\section{Claim \#4}

\paragraph{Claim Statement} ``Theory-of-mind ability has spontaneously emerged in LLMs at a level comparable to that of a 9-year-old child'' \citep{kosinski2024tom}.

\paragraph{Construct Identification} The claim invokes ``theory of mind''---a developmental-psychology construct denoting stable representation of others' mental states (beliefs, desires, intentions, knowledge). The AI Construct Lexis \citep{aiconstructlexis} treats theory of mind as a borrowed cognitive-science construct whose canonical instruments (Sally-Anne, false-belief vignettes) have decades of validation in human children but require independent re-validation for LLMs; the criterion measured is accuracy on Sally-Anne-style false-belief vignettes.

\subsection{Content Validity}
\begin{itemize}
    \item \textbf{Rating:} Partial
    \item \textbf{Confidence:} 3
    \item \textbf{Written Argument:} The Kosinski battery samples one operationalization of theory of mind---false-belief vignettes derived from the Sally-Anne paradigm---but the construct also includes intent recognition, sarcasm and irony comprehension, perspective-taking in dialogue, second-order belief reasoning, and pragmatic inference. Coverage of the construct domain is therefore partial; the evaluation treats accuracy on one narrow sub-task as a proxy for the whole construct. This partial coverage compounds with the deeper portability question raised by \citet{ullman2023llm}: a content sample validated for human children does not automatically transfer to LLMs \citep{salaudeen2025measurement}.
\end{itemize}

\subsection{Construct Validity}
\begin{itemize}
    \item \textbf{Rating:} Weak
    \item \textbf{Confidence:} 5
    \item \textbf{Written Argument:} \textbf{[Selected as the most relevant dimension for this claim.]} Construct validity is the entire game, and \emph{discriminant validity} is the relevant sub-facet. A model with a stable false-belief construct should be invariant to whether a marble is ``in'' or ``inside'' a box, because the construct lives at the level of mental-state attribution rather than lexical surface. \citet{ullman2023llm} shows the opposite: semantically-preserving rewordings of \citet{kosinski2024tom}'s vignettes sharply degrade model performance, the canonical signature that the test is measuring lexical pattern rather than the named construct. \emph{Convergent validity} is also weak in the follow-up literature: models that pass false-belief tasks fail other ToM-adjacent measures (sarcasm, pragmatic reasoning). \emph{Structural validity} is unestablished for the LLM case. Developmental psychologists spent decades showing how performance on false-belief tasks relates to other theory-of-mind sub-skills in children---perspective-taking, deception recognition, sarcasm comprehension. When the same tasks were applied to LLMs, no equivalent mapping work was done. We are using a tool calibrated for one population on a different population without recalibrating it \citep{salaudeen2025measurement}.
\end{itemize}

\subsection{Criterion / Predictive Validity}
\begin{itemize}
    \item \textbf{Rating:} Weak
    \item \textbf{Confidence:} 4
    \item \textbf{Written Argument:} The cited evidence does not supply predictive evidence: no prospective study in the provided materials shows that high false-belief vignette scores predict downstream theory-of-mind-relevant behavior (perspective-taking in dialogue, recognition of communicative intent, audience-appropriate adaptation). Concurrent evidence is also weak: convergent agreement with the broader nomological network of ToM constructs is poor in the follow-up literature, with models passing canonical false-belief tasks failing related social-cognition measures. The criterion evidence (Sally-Anne accuracy) is on a different axis than the construct claim about cross-context mental-state reasoning \citep{ullman2023llm,salaudeen2025measurement}.
\end{itemize}

\subsection{Consequential Validity}
\begin{itemize}
    \item \textbf{Rating:} Weak
    \item \textbf{Confidence:} 4
    \item \textbf{Written Argument:} A ``9-year-old's theory of mind'' claim, taken seriously, justifies deploying LLMs as social agents---therapy chatbots, negotiators, classroom tutors---in which mental-state inference would carry weight. If the construct is illusory but the surface behavior is convincing, the deployment risk is the canonical consequential-validity failure mode: stakeholders confer trust on the basis of measured performance that does not track the relevant underlying capacity \citep{salaudeen2025measurement}. The evidence \citep{ullman2023llm} suggests the underlying construct is not stable, so high-trust deployments invite predictable harm.
\end{itemize}

\subsection{External validity}
\begin{itemize}
    \item \textbf{Rating:} Weak
    \item \textbf{Confidence:} 4
    \item \textbf{Written Argument:} Kosinski's battery uses a narrow set of vignette templates. Generalization across language families, narrative structures, modalities (visual ToM tasks), and task formats is unstudied, and \citet{ullman2023llm} shows that even rewording a vignette without changing its meaning is enough to break the model---generalization fails at the smallest possible scale. The test population (English-language vignettes in a specific developmental tradition) does not stand in for the diverse contexts in which ToM-relevant behavior would matter for any deployed system \citep{salaudeen2025measurement}.
\end{itemize}

\subsection{Overall Validity Judgment}
\begin{itemize}
    \item \textbf{Rating:} Not Valid
    \item \textbf{Confidence:} 4
    \item \textbf{Written Argument:} The score on false-belief vignettes is real and reproducible, but \citet{ullman2023llm} provides direct evidence that the test isn't measuring what it claims to: rewording a vignette without changing its meaning collapses model performance. The model is matching surface word patterns, not reasoning about false beliefs. The ``9-year-old child'' comparison makes things worse---it borrows the framework that, for human children, connects false-belief performance to genuine theory of mind, without doing the work to show that framework applies to LLMs. Claiming human-like developmental milestones on the basis of shortcuts that break under tiny rewordings invites severe consequential risks, especially in social or therapeutic deployments where users expect real mental-state inference. \emph{Defensible reformulation}: ``Several large language models achieve high accuracy on standard false-belief vignettes from developmental psychology. When the same vignettes are reworded without changing their meaning, performance drops sharply \citep{ullman2023llm}, indicating the models are matching surface word patterns rather than reasoning about false beliefs in any stable way. Comparison to human developmental milestones is not supported by the evidence.''
    \item \textbf{Peer Prediction (BTS):}

    \begin{tabular}{ll}
        Valid as stated: & 5\% \\
        Partially valid: & 35\% \\
        Not valid: & 60\% \\
        \end{tabular}\\[4pt]
\end{itemize}
